\subsection{Code Objects, Frame Objects, and Namespace Dictionaries}

After CPython has parsed and compiled source code, the important runtime question becomes: what exactly is executed? The answer is not “the source text.” CPython executes code objects inside frame objects, using namespace dictionaries to resolve names and store bindings.

These three structures form one of the most important runtime triangles in Python:

\begin{verbatim}
code object
    describes what can be executed

frame object
    represents one active execution of that code

namespace dictionaries
    map names to objects during that execution
\end{verbatim}

This distinction is essential because Python code and Python execution are not the same thing. A code object is executable material. A frame is a running activation of that material. A namespace is the name-to-object environment used while that frame executes.

\subsubsection{Code Objects: Stored Executable Descriptions}

A code object is CPython's internal representation of a compiled block of Python code. Module bodies, function bodies, lambdas, comprehensions, class bodies, and nested functions all have code objects.

A code object contains the information needed to execute a block, including:

\begin{itemize}
    \item bytecode instructions;
    \item constants used by the code;
    \item names referenced by the code;
    \item local variable names;
    \item free variables and cell variables for closures;
    \item required evaluation stack size;
    \item filename, function name, and line-number metadata.
\end{itemize}

A simplified view is:

\begin{verbatim}
code object
    bytecode
    constants
    referenced names
    local variable names
    closure variable metadata
    source-location metadata
\end{verbatim}

A code object is not a function call. It does not by itself have current local variable values, an instruction position, or a live execution stack. It is closer to a compiled recipe: it says what instructions exist and what names and constants those instructions refer to.

For example:

\begin{verbatim}
def add(a, b):
    return a + b
\end{verbatim}

When this definition is compiled, the body of \texttt{add} has a code object. That code object knows that there are local names such as \texttt{a} and \texttt{b}, and that the function body contains bytecode for loading them, adding them, and returning the result.

But the code object does not yet know the runtime values of \texttt{a} and \texttt{b}. Those values exist only when the function is called.

\subsubsection{Function Objects Wrap Code Objects}

A function object is not the same thing as a code object. A function object wraps a code object together with runtime context.

Conceptually:

\begin{verbatim}
function object
    code object
    global namespace
    default argument values
    keyword-only defaults
    annotations
    closure cells, if needed
\end{verbatim}

This explains why Python functions are first-class runtime objects. When CPython executes a \texttt{def} statement, it creates a function object and binds the function name to that object.

So this:

\begin{verbatim}
def add(a, b):
    return a + b
\end{verbatim}

does not merely “declare” a function in the C sense. It executes a statement that creates a function object and stores it in the current namespace under the name \texttt{add}.

This is why the following works:

\begin{verbatim}
f = add
result = f(2, 3)
\end{verbatim}

Both \texttt{add} and \texttt{f} refer to the same function object. The function object contains the code object that will be executed when the function is called.

\subsubsection{Frame Objects: One Active Execution}

A frame object represents an active execution of a code object. If a code object is the executable recipe, a frame is one running use of that recipe.

When a function is called, CPython prepares a frame for that call. The frame connects:

\begin{itemize}
    \item the code object being executed;
    \item the current instruction position;
    \item the local names or local slots;
    \item the global namespace;
    \item the built-in namespace;
    \item temporary expression values;
    \item exception and control-flow state.
\end{itemize}

A simplified view is:

\begin{verbatim}
frame
    code object being executed
    current bytecode position
    local namespace / local slots
    global namespace
    builtins namespace
    evaluation stack
    link to caller frame
\end{verbatim}

This is the Python-level equivalent of an activation record, but it should not be confused with a raw C stack frame. CPython is a C program and it uses the native C stack internally, but Python execution also has its own frame model.

This is why Python tracebacks can show a chain of calls. When an exception propagates, CPython can inspect the active or recorded frames and report which code object, line, and call level were involved.

\subsubsection{The Same Code Object Can Have Many Frames}

A code object can be executed many times. Each execution needs its own frame.

For example:

\begin{verbatim}
def add(a, b):
    return a + b

x = add(1, 2)
y = add(10, 20)
\end{verbatim}

The function \texttt{add} has one code object. But the two calls create two separate executions. In the first call, \texttt{a} is bound to \texttt{1} and \texttt{b} to \texttt{2}. In the second call, \texttt{a} is bound to \texttt{10} and \texttt{b} to \texttt{20}.

The code object is shared executable structure. The frame is per-call execution state.

This is the same distinction as:

\begin{verbatim}
recipe       -> code object
cooking once -> frame
\end{verbatim}

The recipe can be reused. Each actual cooking event has its own ingredients, current step, and temporary state.

\subsubsection{Namespace Dictionaries: Names Point to Objects}

Python names live in namespaces. A namespace is a mapping from names to objects. At module level, this mapping is normally a dictionary.

For example:

\begin{verbatim}
x = 10
message = "hello"
\end{verbatim}

creates bindings conceptually like:

\begin{verbatim}
module namespace
    "x"       -> int object 10
    "message" -> str object "hello"
\end{verbatim}

This is why Python variables are not C boxes. A name does not own a fixed typed storage cell. A name is an entry in a namespace, and that entry points to an object.

When Python executes:

\begin{verbatim}
x = 10
x = "hello"
\end{verbatim}

the binding changes:

\begin{verbatim}
before:
    "x" -> int object 10

after:
    "x" -> str object "hello"
\end{verbatim}

The name \texttt{x} has been rebound. It did not mutate from an integer storage box into a string storage box.

\subsubsection{Local, Global, and Built-in Name Resolution}

When code executes, CPython must resolve names. The basic lookup chain for ordinary function code is:

\begin{verbatim}
local namespace
    then global namespace
        then builtins namespace
\end{verbatim}

For example:

\begin{verbatim}
x = 10

def f():
    print(x)
\end{verbatim}

Inside \texttt{f}, the name \texttt{print} is usually found in the built-in namespace. The name \texttt{x} is found in the global namespace, because \texttt{x} is not local to \texttt{f}.

The frame provides access to these namespaces. Conceptually:

\begin{verbatim}
frame for f
    locals   -> local names of f
    globals  -> module dictionary
    builtins -> built-in names dictionary
\end{verbatim}

This explains why Python can resolve a name that was not declared inside the function. It is not reading raw storage from the C stack. It is performing namespace lookup according to Python's execution rules.

\subsubsection{Fast Locals: The Practical Optimization}

Although it is useful to explain namespaces as dictionaries, CPython does not always store ordinary function locals as a normal dictionary during execution. For speed, local variables in optimized function scopes are usually stored internally in a faster array-like layout associated with the frame.

The important conceptual point remains the same: local names are bindings to objects. But the implementation may avoid using a normal dictionary for every local variable access because dictionary lookup would be expensive inside every function call.

So the teaching model is:

\begin{quote}
At the language level, locals behave like a local namespace. At the CPython implementation level, optimized function locals may be stored in faster internal slots and exposed as a mapping when introspection asks for them.
\end{quote}

This explains why \texttt{locals()} inside a function should not be treated as an ordinary mutable dictionary that directly controls execution in all cases. It is a view or mapping of local state, not the same simple thing as a module dictionary.

\subsubsection{Module Namespaces Are Real Dictionaries}

Module namespaces are easier. A module object has a dictionary that stores its top-level names.

For example, a file:

\begin{verbatim}
x = 10

def f():
    return x + 1
\end{verbatim}

creates a module namespace roughly like:

\begin{verbatim}
module.__dict__
    "__name__" -> module name
    "x"        -> int object 10
    "f"        -> function object
\end{verbatim}

The function object \texttt{f} also remembers the global namespace in which it was defined. This is why \texttt{f} can later find \texttt{x}, even if \texttt{f} is passed somewhere else.

For example:

\begin{verbatim}
g = f
\end{verbatim}

does not copy the code or copy the global variables. It binds another name, \texttt{g}, to the same function object. That function object still carries a reference to its original global namespace.

\subsubsection{Class Body Execution Uses a Namespace Too}

Class definitions also use namespaces, but they behave differently from ordinary functions.

When Python executes:

\begin{verbatim}
class Point:
    x = 0
    y = 0

    def move(self, dx, dy):
        self.x += dx
        self.y += dy
\end{verbatim}

the class body is executed in a temporary namespace. Names assigned in the class body become entries in that namespace. After the body finishes, Python creates a class object using that namespace.

Conceptually:

\begin{verbatim}
temporary class namespace
    "x"    -> int object 0
    "y"    -> int object 0
    "move" -> function object

type machinery creates class object Point
Point.__dict__ contains these entries
\end{verbatim}

This is why a class definition is executable code, not only a static declaration. The body runs, creates objects, binds names, and then those names become the class object's attributes.

\subsubsection{Closures: When Frames Need Preserved Cells}

Nested functions introduce one more mechanism: closure cells.

Example:

\begin{verbatim}
def outer():
    x = 10

    def inner():
        return x

    return inner
\end{verbatim}

Here, \texttt{inner} uses \texttt{x} from \texttt{outer}. But \texttt{outer} returns, so its ordinary call frame cannot simply remain active as a normal stack frame forever. CPython solves this by storing captured variables in cell objects. The returned function keeps references to those cells.

Conceptually:

\begin{verbatim}
inner function object
    code object
    globals
    closure cells
        cell for x -> int object 10
\end{verbatim}

This is the concrete runtime basis of closures. A closure is not a mystical memory of the past. It is a function object carrying references to cells that preserve needed bindings from an enclosing scope.

\subsubsection{Why This Matters for Python Programming}

This machinery explains several everyday Python facts.

First, function definitions create objects:

\begin{verbatim}
def f():
    pass
\end{verbatim}

After this executes, \texttt{f} is a name bound to a function object.

Second, assignment binds names:

\begin{verbatim}
a = []
b = a
\end{verbatim}

Both names point to the same list object.

Third, function calls create frames:

\begin{verbatim}
f()
\end{verbatim}

Calling \texttt{f} does not copy the source text. It creates execution state for \texttt{f}'s code object.

Fourth, globals are namespace entries:

\begin{verbatim}
x = 5
\end{verbatim}

At module level, this stores \texttt{x} in the module dictionary.

Fifth, closures preserve selected bindings:

\begin{verbatim}
def outer():
    x = 5
    def inner():
        return x
    return inner
\end{verbatim}

The returned \texttt{inner} function can still access \texttt{x} because the needed binding is preserved in closure machinery.

\subsubsection{The Mental Model}

The complete model is:

\begin{verbatim}
source code
    is compiled into code objects

function definitions
    create function objects that wrap code objects

function calls
    create frames that execute code objects

frames
    contain or access local, global, and built-in namespaces

namespaces
    bind names to objects

objects
    carry identity, type, value, and implementation metadata
\end{verbatim}

This is the correct level of detail for learning Python effectively. We do not need to memorize CPython's source files, but we must understand the architecture.

Python execution is not text being read casually. It is compiled code objects running inside frame state, resolving names through namespace mappings, and manipulating objects through the CPython runtime.